% Source: Weinberg GR, physical PDF pp. 243--245
%         (printed pp. 220--222).
% Coverage: Section 9.2, Particle and Photon Dynamics; equations
%           (9.2.1)--(9.2.7).
% Figures/tables/footnotes: no figures or tables.
% Status: source-reviewed and compile-clean.
% Uncertainties: none.

\subsection{Particle and Photon Dynamics}\label{sec:9.2}

Before continuing our calculation of the post-Newtonian metric, we shall
take a quick look back at the problem with which we started, that of
computing the acceleration of a freely falling particle to order
\(\bar v^4/\bar r\). (Detailed applications of the post-Newtonian method
are given in Sections~\ref{sec:9.5}--\ref{sec:9.9}.) Inserting
Eqs.~\eqref{eq:9.1.67}--\eqref{eq:9.1.72} into Eq.~\eqref{eq:9.1.2}
immediately gives the equation of motion
\begin{equation}
  \begin{split}
  \frac{\dd\mathbf v}{\dd t}={}&
  -\boldsymbol\nabla(\phi+2\phi^2+\psi)
  -\partial_t\boldsymbol\zeta
  +\mathbf v\times(\boldsymbol\nabla\times\boldsymbol\zeta)\\
  &+3\mathbf v\,\partial_t\phi
  +4\mathbf v(\mathbf v\cdot\boldsymbol\nabla\phi)
  -v^2\boldsymbol\nabla\phi,
  \qquad
  \mathbf v\equiv\frac{\dd\mathbf x}{\dd t}.
  \end{split}
  \tag{9.2.1}\label{eq:9.2.1}
\end{equation}

In addition, we need to know how to convert the harmonic coordinate time
\(t\) into the proper time \(\tau\) measured on a freely falling body of
velocity \(\mathbf v\). By definition,
\[
  \left(\frac{\dd\tau}{\dd t}\right)^2
  =-g_{00}-2g_{i0}v^i-g_{ij}v^iv^j.
\]
To order \(\bar v^4\), this is
\[
  \left(\frac{\dd\tau}{\dd t}\right)^2
  =
  1-\bigl[v^2+g_{00}^{(2)}\bigr]
  -\bigl[g_{00}^{(4)}+2g_{i0}^{(3)}v^i
          +g_{ij}^{(2)}v^iv^j\bigr],
\]
or, using Eqs.~\eqref{eq:9.1.57}, \eqref{eq:9.1.60},
\eqref{eq:9.1.61}, and \eqref{eq:9.1.63},
\[
  \left(\frac{\dd\tau}{\dd t}\right)^2
  =
  1+(2\phi-v^2)
  +2(\phi^2+\psi-\boldsymbol\zeta\cdot\mathbf v+\phi v^2).
\]
The parentheses enclose terms of orders \(\bar v^2\) and \(\bar v^4\).
Expanding the square root gives
\[
  \frac{\dd\tau}{\dd t}
  =
  1+\phi-\frac12v^2-\frac18(2\phi-v^2)^2
  +\phi^2+\psi-\boldsymbol\zeta\cdot\mathbf v+\phi v^2,
\]
or
\begin{equation}
  \frac{\dd\tau}{\dd t}=1-L,
  \tag{9.2.2}\label{eq:9.2.2}
\end{equation}
where
\begin{equation}
  L
  =
  \frac12v^2-\phi-\frac12\phi^2-\frac32\phi v^2
  +\frac18(v^2)^2-\psi+\boldsymbol\zeta\cdot\mathbf v.
  \tag{9.2.3}\label{eq:9.2.3}
\end{equation}
Since \(\int(\dd\tau/\dd t)\,\dd t\) is stationary, we can regard \(L\) as
the Lagrangian for a single particle and derive its equations of motion from
\begin{equation}
  \frac{\dd}{\dd t}\frac{\partial L}{\partial v^i}
  =
  \frac{\partial L}{\partial x^i}.
  \tag{9.2.4}\label{eq:9.2.4}
\end{equation}
(Acting on \(\phi\) or \(\boldsymbol\zeta\),
\(\dd/\dd t\) is to be taken as
\(\partial_t+\mathbf v\cdot\boldsymbol\nabla\).) The reader may readily
check that Eq.~\eqref{eq:9.2.4} agrees with Eq.~\eqref{eq:9.2.1}.

The post-Newtonian fields can also be used to calculate the acceleration of
a photon in a gravitational field to order \(\bar v^2\). (Here \(\bar v\)
is of course not the photon speed; it is the typical velocity of the
material particles of the system.) Since the velocity
\(\mathbf u\equiv\dd\mathbf x/\dd t\) of the photon is of order unity,
Eq.~\eqref{eq:9.1.2} gives
\[
  \frac{\dd u_i}{\dd t}
  =
  -\Gamma^{i(2)}{}_{00}
  -\Gamma^{i(2)}{}_{jk}u_ju_k
  +2u_i\Gamma^{0(2)}{}_{0j}u_j
  +\order(\bar v^3).
\]
Using Eqs.~\eqref{eq:9.1.67}, \eqref{eq:9.1.70}, and
\eqref{eq:9.1.72}, this gives
\[
  \frac{\dd\mathbf u}{\dd t}
  =
  -(1+u^2)\boldsymbol\nabla\phi
  +4\mathbf u(\mathbf u\cdot\boldsymbol\nabla\phi)
  +\order(\bar v^3).
\]
The photon speed follows from
\[
  0
  =
  -g_{\mu\nu}\frac{\dd x^\mu}{\dd t}\frac{\dd x^\nu}{\dd t}
  =
  1-u^2+2(1+u^2)\phi+\order(\bar v^3),
\]
so
\begin{equation}
  |\mathbf u|=1+2\phi+\order(\bar v^3).
  \tag{9.2.5}\label{eq:9.2.5}
\end{equation}
Hence, to the required accuracy, we may replace \(u^2\) by unity in the
photon acceleration:
\begin{equation}
  \frac{\dd\mathbf u}{\dd t}
  =
  -2\boldsymbol\nabla\phi
  +4\mathbf u(\mathbf u\cdot\boldsymbol\nabla\phi)
  +\order(\bar v^3).
  \tag{9.2.6}\label{eq:9.2.6}
\end{equation}
It is somewhat more convenient to write this as an equation for the unit
direction vector \(\widehat{\mathbf u}\equiv\mathbf u/|\mathbf u|\):
\begin{equation}
  \frac{\dd\widehat{\mathbf u}}{\dd t}
  =
  \widehat{\mathbf u}\times
  (\widehat{\mathbf u}\times\boldsymbol\nabla\phi)
  +\order(\bar v^3).
  \tag{9.2.7}\label{eq:9.2.7}
\end{equation}
